\documentclass[11pt]{article}

\usepackage[a4paper,margin=27mm]{geometry}
\usepackage[T1]{fontenc}
\usepackage[utf8]{inputenc}
\usepackage{lmodern}
\usepackage{microtype}
\usepackage{amsmath,amssymb,amsthm,mathtools}
\usepackage{booktabs,longtable,array,tabularx}
\usepackage{enumitem}
\usepackage{xcolor}
\usepackage{hyperref}
\usepackage{cleveref}
\usepackage{listings}
\usepackage{fancyhdr}
\usepackage{lastpage}
\usepackage{needspace}
\usepackage[most]{tcolorbox}

\hypersetup{
  colorlinks=true,
  linkcolor=blue!55!black,
  citecolor=green!35!black,
  urlcolor=blue!60!black,
  pdftitle={The Alaoglu--Erdos Conjecture at Weight Two},
  pdfauthor={OpenAI, prepared for Vladimir Reshetnikov}
}

\setlist{nosep,leftmargin=2em}
\setlength{\parindent}{0pt}
\setlength{\parskip}{0.55em}
\setlength{\emergencystretch}{3em}
\setlength{\headheight}{14pt}

\pagestyle{fancy}
\fancyhf{}
\fancyhead[L]{Alaoglu--Erd\H{o}s at weight two}
\fancyhead[R]{August 21, 2026}
\fancyfoot[C]{\thepage\ of \pageref{LastPage}}

\newtheorem{theorem}{Theorem}[section]
\newtheorem{proposition}[theorem]{Proposition}
\newtheorem{lemma}[theorem]{Lemma}
\newtheorem{corollary}[theorem]{Corollary}
\newtheorem{conjecture}[theorem]{Conjecture}
\theoremstyle{definition}
\newtheorem{definition}[theorem]{Definition}
\newtheorem{remark}[theorem]{Remark}
\newtheorem{warning}[theorem]{Audit warning}

\definecolor{provedbg}{RGB}{235,247,238}
\definecolor{compbg}{RGB}{238,244,252}
\definecolor{openbg}{RGB}{253,242,234}
\newtcolorbox{provedbox}{
  enhanced,breakable,colback=provedbg,colframe=green!35!black,
  boxrule=0.55pt,arc=1.2mm,left=2mm,right=2mm,top=1.5mm,bottom=1.5mm,
  before upper={\textbf{PROVED.}\par\smallskip}
}
\newtcolorbox{compbox}{
  enhanced,breakable,colback=compbg,colframe=blue!35!black,
  boxrule=0.55pt,arc=1.2mm,left=2mm,right=2mm,top=1.5mm,bottom=1.5mm,
  before upper={\textbf{COMPUTATION.}\par\smallskip}
}
\newtcolorbox{conditionalbox}{
  enhanced,breakable,colback=openbg,colframe=orange!55!black,
  boxrule=0.55pt,arc=1.2mm,left=2mm,right=2mm,top=1.5mm,bottom=1.5mm,
  before upper={\textbf{OPEN / CONDITIONAL.}\par\smallskip}
}

\newcommand{\N}{\mathbb N}
\newcommand{\Z}{\mathbb Z}
\newcommand{\Q}{\mathbb Q}
\newcommand{\R}{\mathbb R}
\newcommand{\C}{\mathbb C}
\newcommand{\Li}{\operatorname{Li}}
\newcommand{\supp}{\operatorname{supp}}
\newcommand{\rank}{\operatorname{rank}}
\newcommand{\Span}{\operatorname{span}}
\newcommand{\val}{\boldsymbol\nu}
\newcommand{\height}{h}
\newcommand{\D}{\mathcal D}
\newcommand{\Kc}{\mathcal C}
\newcommand{\Vdef}{\mathcal V}
\newcommand{\odotQ}{\mathbin{\odot}}
\newcommand{\sha}[1]{\texttt{#1}}

\lstdefinestyle{shell}{
  basicstyle=\ttfamily\small,
  columns=fullflexible,
  breaklines=true,
  frame=single,
  backgroundcolor=\color{black!2},
  showstringspaces=false
}

\title{\textbf{The Alaoglu--Erd\H{o}s Conjecture at Weight Two}\\[0.35em]
\large A Pad\'e Height Audit, Rational-Transfer No-Go Theorems,\\
and a Directed-Rounding Extension through $2^{29}$}
\author{Prepared for Vladimir Reshetnikov\\
\small A nonduplicative continuation of the supplied unified report}
\date{August 21, 2026}

\begin{document}
\maketitle

\begin{abstract}
Let $x\in\R$ and suppose that $2^x$ and $3^x$ are integers.  The Alaoglu--Erd\H{o}s
conjecture asserts that $x$ is an integer.  The supplied 327-page unified report already
reduces a hypothetical counterexample to an exact rank-two monoid, develops many determinant,
valuation, dynamical, matrix, and operator formulations, and identifies the remaining wall as
an unproved weight-two relation between logarithms.  This report does not repeat that material.

The full conjecture is not proved here.  Instead, the report gives a new and fairly sharp audit
of the most relevant post-report development: Kawashima's May 2026 Pad\'e criterion for products
of polylogarithms.  The central new theorem is a representation-independent upper bound for
Kawashima's height defect.  It shows that a successful weight-two specialization requires the
polylogarithm arguments to have almost maximal \emph{arithmetic efficiency}: their local
smallness must be within a factor $1-1/(m+1)^2$ of their average global height.  This one
inequality rigorously rules out the direct prime-logarithm points, high-root splitting,
ordinary telescoping decompositions, and every unbounded construction inside a fixed
$S$-unit group.

A second new theorem goes further.  At rational points, Kawashima's criterion is incompatible
with having even $\log 2$ in the rational span of the resulting logarithms.  The proof combines
the height-defect bound with a Cramer/Hadamard bound on a prime-valuation matrix.  A tensor-rank
corollary rules out the corresponding formal quadratic transfer of the entire
Alaoglu--Erd\H{o}s relation.  Thus the current Pad\'e theorem cannot settle the conjecture by
any rational-point encoding of the natural kind; a successful continuation must exploit a
genuinely new mixed-denominator or mixed-embedding cancellation mechanism.

Finally, the exact MPFR verifier reproduced in the supplied report was rerun, without source
changes, on the two new bit ranges $[2^{27},2^{28})$ and $[2^{28},2^{29})$.  All
$402{,}653{,}182$ non-powers of two were separated from the nearest integer by directed
logarithmic inequalities.  At the stated software trust level this extends the finite
exclusion to $2^x<2^{29}$, so any nonintegral counterexample must satisfy $x>29$.
\end{abstract}

\tableofcontents
\newpage

\section{Executive verdict and claim ledger}

\begin{conditionalbox}
The requested unconditional proof was not obtained.  As of the literature check performed for
this report, the problem remains an open special case of the Four Exponentials Conjecture
\cite{MOAE}.  No
step below is promoted from conditional or computational status to theorem status merely to
produce a desired conclusion.
\end{conditionalbox}

The work nevertheless changes the technical picture in three concrete ways.

\begin{enumerate}[label=\textbf{\arabic*.}]
\item It extracts from Kawashima's criterion an intrinsic, representation-free obstruction
      (\Cref{thm:projective-defect}) and turns it into an arithmetic-efficiency threshold.
\item It proves that the entire rational-point transfer strategy is impossible under that
      criterion (\Cref{thm:rational-transfer,cor:quadratic-transfer}), not merely that several
      obvious choices of a common denominator fail.
\item It extends the independent finite scan by two full bits, from $2^{27}$ to $2^{29}$
      (\Cref{sec:scan}).
\end{enumerate}

\begin{center}
\begin{tabularx}{\textwidth}{@{}>{\raggedright\arraybackslash}p{0.28\textwidth}
                                  >{\raggedright\arraybackslash}p{0.16\textwidth}X@{}}
\toprule
Statement & Status & Logical role \\
\midrule
Projective Pad\'e defect bound & Proved & Necessary condition for every specialization of Kawashima's weight-two criterion. \\
Arithmetic-efficiency threshold & Proved & Quantifies the exact high-efficiency region reached by the present theorem. \\
Direct, radical, and telescoping encodings fail & Proved & Eliminates the most natural representations of integer logarithms. \\
Rational-span transfer impossibility & Proved & Shows that no rational specialization satisfying the criterion can even span $\log 2$. \\
Formal quadratic-transfer impossibility & Proved & Extends the preceding no-go from factor-by-factor substitutions to the natural prime-valuation tensor of a counterexample. \\
Fixed-$S$ exceptional finiteness & Proved from published effective $S$-unit approximation & Any fixed multiplicative group supplies at most effectively many candidate Pad\'e tuples. \\
Near-unit unimodular form & Proved & Converts a counterexample into infinitely many exact determinant zeros near $1$, but in the low-efficiency regime. \\
Scan through $2^{29}$ & Directed-rounding computation & Excludes nonintegral $x\le 29$ at a disclosed software trust boundary. \\
Full Alaoglu--Erd\H{o}s conjecture & Open & Still requires a genuinely new weight-two transcendence theorem. \\
\bottomrule
\end{tabularx}
\end{center}

\section{Scope: what is imported and what is not repeated}

Write
\[
  M=2^x,\qquad A=3^x.
\]
The supplied unified report is taken as the source for the following previously established
facts and constructions:

\begin{itemize}
\item rational $x$ is automatically integral, while a nonintegral solution is transcendental;
\item a counterexample gives the exact determinant relation
\begin{equation}\label{eq:wall}
   \log M\,\log 3-\log A\,\log 2=0;
\end{equation}
\item conditional on one counterexample, all solutions form the free additive monoid
      $\{n+k\beta:n,k\ge0\}$ for a unique least nonintegral generator $\beta$;
\item the corresponding output pairs form the free multiplicative monoid generated by
      $(2,3)$ and $(M,A)$;
\item the ordinary Four Exponentials Conjecture implies the desired conclusion, whereas the
      Six Exponentials Theorem does not close the two-by-two case;
\item the two outputs of a nonintegral solution cannot both be $3$-smooth;
\item the existing directed-rounding audit covers $2^x<2^{27}$.
\end{itemize}

The long determinant, valuation, Sturmian, matrix, operator, and formalization developments of
the supplied report are not copied.  Only the minimal statements needed to connect the new
Pad\'e analysis to the original problem are restated.

\section{The exact weight-two object}

\subsection{Prime-valuation form}

Let $\mathcal P$ be the set of rational primes and let
\[
  E=\Q^{(\mathcal P)}
\]
be the finite-support rational vector space with basis $(e_p)_{p\in\mathcal P}$.  For a
positive rational number $r$, write
\[
  \val(r)=(v_p(r))_{p\in\mathcal P}\in E,
  \qquad
  \lambda(y)=\sum_p y_p\log p.
\]
Unique factorization gives an elementary but important fact:
\begin{equation}\label{eq:lambda-inj}
  \lambda:E\longrightarrow\R\quad\text{is injective as a $\Q$-linear map.}
\end{equation}
Indeed, after clearing denominators, $\lambda(y)=0$ says that a positive rational number is
$1$.

Put
\[
  u=\val(M),\qquad v=\val(A),\qquad a=e_2,\qquad b=e_3.
\]
Then \eqref{eq:wall} is the vanishing at the point $\lambda$ of the rational symmetric tensor
\begin{equation}\label{eq:tensor}
  T=u\odotQ b-v\odotQ a\in\operatorname{Sym}^2(E),
  \qquad y\odotQ z:=\frac{y\otimes z+z\otimes y}{2}.
\end{equation}
This is the exact weight-two object to which a product-of-logarithms theorem must apply.

\subsection{Rank of the counterexample tensor}

The following elementary rank audit will be used later to prevent a hidden low-dimensional
escape in a formal Pad\'e transfer.

\begin{lemma}[Full range on the support span]\label{lem:tensor-rank}
Assume that $(M,A)$ comes from a nonintegral counterexample.  Let
\[
  S=\Span_\Q\{u,v,a,b\}.
\]
Then $\dim S\ge3$, the tensor $T$ of \eqref{eq:tensor} has rank $\dim S$, and consequently
its contraction range is exactly $S$.
\end{lemma}

\begin{proof}
If $\dim S=2$, then $u$ and $v$ lie in $\Span\{e_2,e_3\}$, so both $M$ and $A$ are
$3$-smooth.  This is excluded by the three-base/Six Exponentials consequence proved in the
supplied report.

If $\dim S=4$, use the four independent linear forms $u,b,v,a$ as coordinates on $S^*$.
The quadratic form associated with $T$ is
\[
  X_1X_2-X_3X_4,
\]
whose matrix is the direct sum of two nondegenerate hyperbolic planes.  Its rank is $4$.

Suppose $\dim S=3$.  The rank is at most $3$, and $T\ne0$: equality $u\odotQ b=v\odotQ a$
would force the two supporting planes to coincide, hence $\dim S\le2$.  If the rank were $1$,
the rational quadratic form would have the shape $c\,\xi(r)^2$ with $c\in\Q^\times$ and
$0\ne r\in S$.  Evaluating at $\lambda$ and using \eqref{eq:wall} would give
$\lambda(r)=0$, contrary to \eqref{eq:lambda-inj}.

It remains to exclude rank $2$.  The vectors $u,v$ are independent: otherwise
$\log M/\log A$ would be rational, whereas it equals $\log2/\log3$.  The vectors $a,b$ are
also independent.  Hence the annihilators of $\Span\{u,v\}$ and $\Span\{a,b\}$ are two
distinct rational isotropic lines.  The radical of a rank-two form on a three-dimensional
space is only one-dimensional, so at least one of these isotropic lines is nonradical.  The
rank-two quotient is therefore a rational hyperbolic plane, and the form factors as
\[
  q_T(\xi)=c\,\xi(r)\xi(s)
\]
with $c\in\Q^\times$ and nonzero $r,s\in S$.  Evaluating at $\lambda$ gives
$\lambda(r)\lambda(s)=0$.  Injectivity \eqref{eq:lambda-inj} again yields a contradiction.
Thus the rank is $3$.

In either possible dimension, rank equals $\dim S$, so the range is all of $S$.
\end{proof}

\section{The 2024--2026 frontier}

Three recent developments are especially close to the remaining wall.

\paragraph{Matrix coefficients.}
Dasgupta and Kakde formulate the Matrix Coefficient Conjecture for matrices of complex or
$p$-adic logarithms.  In size two it contains the Four Exponentials problem rather than
solving it: a zero determinant should become a zero coefficient after rational changes of
basis.  Their work clarifies the algebraic shape of the obstruction, but the required
coefficient theorem remains conjectural in the present case \cite{DasguptaKakde2026}.

\paragraph{Arithmetic holonomy.}
Calegari, Dimitrov, and Tang prove strong arithmetic-holonomy bounds and, among other
applications, prove the $\Q$-linear independence of
\[
  1,\quad \log\!\left(1+\frac1m\right),\quad
  \log\!\left(1+\frac1n\right),\quad
  \log\!\left(1+\frac1m\right)\log\!\left(1+\frac1n\right)
\]
when $m/n$ is sufficiently close to $1$ and $m\ne n$ \cite{CDTProducts}.  The same authors'
effective approximation theorem gives a uniform height bound for solutions of
$|1-A\gamma|_v\le H(\gamma)^{-\varepsilon}$ in a fixed finitely generated multiplicative
group \cite{CDTHolonomy}.  These are real advances at weight two, but their proven product
case lies on the special unit-gap locus $(m+1)/m$.

\paragraph{Products of polylogarithms.}
Kawashima proves a new Pad\'e criterion for linear independence of products of
polylogarithms at distinct algebraic points \cite{Kawashima2026}.  Since
\[
  \Li_1(z)=-\log(1-z),
\]
this is the first theorem in the new literature whose conclusion directly contains mixed
products of logarithms at many points.  The rest of this report audits exactly how far its
height hypothesis reaches toward \eqref{eq:wall}.

\section{Kawashima's weight-two defect}

Let $K$ be a number field, $v$ a place, $m\ge1$, and
$\boldsymbol\alpha=(\alpha_1,\ldots,\alpha_m)\in(K^\times)^m$ with pairwise distinct
coordinates.  Put
\[
  \D_m=(m+1)^2-1=m^2+2m
\]
and
\begin{equation}\label{eq:Cm}
  \Kc_m=\D_m\log2+3\log(m+1)+2+2\D_m.
\end{equation}
For $\beta\in K$ satisfying $|\beta|_v>H_v(\boldsymbol\alpha)$, Kawashima's weight-two defect
is
\begin{align}\label{eq:Kdefect}
 \Vdef_v(\boldsymbol\alpha,\beta)
  ={}&(\D_m+1)h_v(\beta)-h_v(\boldsymbol\alpha)\\
    &-\D_m\left(h(\beta)+\frac1m\sum_{i=1}^m h(\alpha_i)
                     +h(\boldsymbol\alpha)\right)-\Kc_m.\notag
\end{align}
If this number is positive, then $1$ together with all products of polylogarithms of total
weight at most two at the points $\alpha_i/\beta$ is linearly independent over $K$.
In particular, all numbers
\[
  1,\qquad \Li_1(\alpha_i/\beta),\qquad
  \Li_1(\alpha_i/\beta)\Li_1(\alpha_j/\beta)
\]
are linearly independent over $K$.

The formula \eqref{eq:Kdefect} looks representation-dependent: the same points
$z_i=\alpha_i/\beta$ admit many common scalings.  The next theorem removes that ambiguity by
bounding the defect only in terms of the points themselves.

\section{A representation-independent Pad\'e obstruction}

\subsection{The projective defect theorem}

\begin{theorem}[Projective height-defect bound]\label{thm:projective-defect}
Keep the notation above and put
\[
  z_i=\frac{\alpha_i}{\beta},\qquad
  a_v(\mathbf z)=-\log\max_i|z_i|_v,
  \qquad
  \overline h(\mathbf z)=\frac1m\sum_{i=1}^m h(z_i).
\]
Then $a_v(\mathbf z)>0$ and
\begin{equation}\label{eq:projective-bound}
  \boxed{\quad
  \Vdef_v(\boldsymbol\alpha,\beta)
  \le (\D_m+1)a_v(\mathbf z)-\D_m\overline h(\mathbf z)-\Kc_m.
  \quad}
\end{equation}
The right-hand side is independent of the choice of common numerator and denominator.
\end{theorem}

\begin{proof}
Write
\[
 t=h_v(\beta)=\log|\beta|_v,
 \qquad A=h_v(\boldsymbol\alpha).
\]
The strict size hypothesis implies $|\beta|_v>1$.  Since
$\max_i|\alpha_i|_v=|\beta|_v\max_i|z_i|_v$, we have
\begin{equation}\label{eq:Ata}
  A=\max\{0,t-a_v(\mathbf z)\}.
\end{equation}
Set
\[
 H=h(\beta),\qquad X=\frac1m\sum_i h(\alpha_i),\qquad
 Y=h(\boldsymbol\alpha),\qquad P=\overline h(\mathbf z).
\]
The elementary height inequality $h(\alpha_i/\beta)\le h(\alpha_i)+h(\beta)$ gives
$P\le X+H$.  Also $H\ge t$ and $Y\ge A$.  Therefore
\begin{equation}\label{eq:HXY}
  H+X+Y\ge \max\{t,P\}+A.
\end{equation}
Substituting \eqref{eq:HXY} into \eqref{eq:Kdefect} yields
\begin{equation}\label{eq:piecewise}
 \Vdef_v\le
 (\D_m+1)t-(\D_m+1)\max\{0,t-a_v\}
 -\D_m\max\{t,P\}-\Kc_m.
\end{equation}
Finally,
\[
 h(z_i)=h(z_i^{-1})\ge h_v(z_i^{-1})=-\log|z_i|_v,
\]
so $P\ge a_v$.  The piecewise-linear expression in \eqref{eq:piecewise} is then maximized
for any $t$ in the interval between $a_v$ and $P$, where its value is
$(\D_m+1)a_v-\D_mP-\Kc_m$.  This proves \eqref{eq:projective-bound}.
\end{proof}

\begin{provedbox}
\Cref{thm:projective-defect} is the main new analytic theorem of this report.  It applies at
archimedean and nonarchimedean places, over arbitrary number fields, and optimizes the local scaling variable in the universal projective bound over every
common scaling of $(\boldsymbol\alpha,\beta)$.
\end{provedbox}

\subsection{Arithmetic efficiency}

\begin{definition}
For a tuple with $|z_i|_v<1$, define its arithmetic efficiency at $v$ by
\[
  \eta_v(\mathbf z)=\frac{a_v(\mathbf z)}{\overline h(\mathbf z)}\in(0,1].
\]
\end{definition}

\begin{corollary}[Efficiency and depth thresholds]\label{cor:efficiency}
If $\Vdef_v(\boldsymbol\alpha,\beta)>0$, then
\begin{align}
  a_v(\mathbf z)&>\Kc_m,\label{eq:depth}\\
  \eta_v(\mathbf z)&>\frac{\D_m}{\D_m+1}
      =1-\frac1{(m+1)^2},\label{eq:eta}\\
  \overline h(\mathbf z)&<\left(1+\frac1{\D_m}\right)a_v(\mathbf z).
  \label{eq:Pupper}
\end{align}
\end{corollary}

\begin{proof}
Since $\overline h\ge a_v$, positivity in \eqref{eq:projective-bound} first gives
$a_v>\Kc_m$.  Dropping the positive constant from the same inequality gives
$(\D_m+1)a_v>\D_m\overline h$, which is \eqref{eq:eta}; rearrangement gives
\eqref{eq:Pupper}.
\end{proof}

The numerical severity is visible immediately.

\begin{center}
\begin{tabular}{@{}rrrrr@{}}
\toprule
$m$ & $\D_m$ & $\Kc_m$ & $e^{-\Kc_m}$ & required $\eta_v$ \\
\midrule
1 & 3  & 12.158883 & $5.24\times10^{-6}$  & $>0.750000$ \\
2 & 8  & 26.841014 & $2.20\times10^{-12}$ & $>0.888889$ \\
3 & 15 & 46.556091 & $6.04\times10^{-21}$ & $>0.937500$ \\
4 & 24 & 71.463846 & $9.20\times10^{-32}$ & $>0.960000$ \\
5 & 35 & 101.635430& $7.25\times10^{-45}$ & $>0.972222$ \\
6 & 48 & 137.108795& $2.85\times10^{-60}$ & $>0.979592$ \\
\bottomrule
\end{tabular}
\end{center}

Thus ``make all points small'' is not enough.  Their local smallness must consume almost all
of their global arithmetic height.

\subsection{Exact scaling loss for integral presentations}

\begin{proposition}[Common scaling is strictly harmful]\label{prop:scaling}
Let $K=\Q$, $v=\infty$, and let $1\le\alpha_i<\beta$ be positive integers.  Put
$\alpha_{\max}=\max_i\alpha_i$.  Then
\begin{equation}\label{eq:integer-defect}
 \Vdef_\infty(\boldsymbol\alpha,\beta)
 =\log\beta-(\D_m+1)\log\alpha_{\max}
   -\frac{\D_m}{m}\sum_i\log\alpha_i-\Kc_m.
\end{equation}
For every integer $t\ge1$,
\begin{equation}\label{eq:scaling-law}
 \Vdef_\infty(t\boldsymbol\alpha,t\beta)
 =\Vdef_\infty(\boldsymbol\alpha,\beta)-2\D_m\log t.
\end{equation}
\end{proposition}

\begin{proof}
For positive integers, every finite-place local height vanishes, so
$h(\beta)=h_\infty(\beta)=\log\beta$,
$h(\alpha_i)=\log\alpha_i$, and
$h(\boldsymbol\alpha)=h_\infty(\boldsymbol\alpha)=\log\alpha_{\max}$.
Substitution gives \eqref{eq:integer-defect}; scaling all terms gives
\eqref{eq:scaling-law}.
\end{proof}

This exact law explains why clearing a larger common denominator cannot rescue a failed
specialization.  It also corrects a tempting but false intuition: because the polylogarithm
arguments are unchanged, one might expect the criterion to be projectively neutral.  It is
not; the arithmetic approximants remember the chosen affine presentation.

\section{Natural encodings and why each misses the criterion}

\subsection{Direct prime-logarithm points}

For every integer $q\ge2$,
\begin{equation}\label{eq:direct}
  \log q=\Li_1\!\left(1-\frac1q\right).
\end{equation}
But $1-1/q\ge1/2$.  Hence any tuple containing a direct point has
$a_\infty\le\log2$, while \Cref{cor:efficiency} requires
$a_\infty>\Kc_m\ge\Kc_1>12$.  Therefore:

\begin{corollary}[Direct points never qualify]\label{cor:direct-no}
No specialization of Kawashima's weight-two criterion containing even one point
$1-1/q$ with $q\ge2$ can have positive defect, regardless of the common representation or
of the other points.
\end{corollary}

This rules out the immediate choice
\[
  1/2,\quad 2/3,\quad 1-1/M,\quad 1-1/A
\]
for the four logarithms in \eqref{eq:wall}.

\subsection{High-root splitting and degree dilution}

A natural attempt to make the direct point small is
\begin{equation}\label{eq:root-split}
  \log q=n\Li_1(1-q^{-1/n}).
\end{equation}
For $q=2$ or $3$, the polynomial $X^n-q$ is Eisenstein, so every number field containing
$q^{1/n}$ has degree at least $n$.  At the positive real place, Kawashima's normalized
absolute value therefore gives
\[
 a_v=-\frac1{[K:\Q]}\log(1-q^{-1/n})
 \le \frac1n\log\!\left(1+\frac{n}{\log q}\right)<1.
\]
Here we used $1-e^{-y}\ge y/(1+y)$ for $y>0$.  The depth condition
$a_v>\Kc_m$ is impossible.

\begin{corollary}[Radical splitting never qualifies]\label{cor:radical-no}
For either base logarithm $\log2$ or $\log3$, replacing it by arbitrary high-root identities
of the form \eqref{eq:root-split} cannot satisfy Kawashima's criterion at a real place.  The
apparent archimedean gain is more than canceled by the normalization factor $[K:\Q]^{-1}$.
\end{corollary}

This is a useful warning for algebraic-point constructions: increasing the degree to make one
chosen conjugate spectacularly close to $1$ can reduce the favorable local height by exactly
the same degree.

\subsection{Telescoping by unit fractions}

For integers $q\ge2$ and $N\ge1$,
\begin{equation}\label{eq:telescoping}
 \log q=\sum_{k=N+1}^{qN}\log\frac{k}{k-1}
       =\sum_{k=N+1}^{qN}\Li_1\!\left(\frac1k\right).
\end{equation}
These points have perfect individual efficiency: for $z=1/k$,
$-\log z=h(z)=\log k$.  Nevertheless the number of distinct points is
$m=(q-1)N$, while the largest point is $1/(N+1)$.  Thus
\[
  a_\infty=\log(N+1),\qquad
  \Kc_m>2m^2.
\]
Since $\log(N+1)<2(q-1)^2N^2$ for every $q\ge2,N\ge1$, the necessary depth condition
$a_\infty>\Kc_m$ fails before any finer height calculation is needed.

\begin{corollary}[Dimension tax for ordinary telescoping]\label{cor:telescoping-no}
No ordinary telescoping decomposition \eqref{eq:telescoping} satisfies Kawashima's
weight-two criterion.
\end{corollary}

If one nevertheless clears all denominators with
$L=\operatorname{lcm}(N+1,\ldots,qN)$, \Cref{prop:scaling} makes the failure substantially
worse.  The exact integral defect is
\[
 -2\D_m\log L+(\D_m+1)\log(N+1)
 +\frac{\D_m}{m}\log\frac{(qN)!}{N!}-\Kc_m,
\]
which tends to $-\infty$ rapidly.  The projective theorem shows that this is not a defect of
that particular clearing; no other common scaling can reverse the sign.

\section{A universal no-go theorem for rational Pad\'e transfer}

The preceding examples might suggest searching for a more ingenious rational family.  The
next theorem shows that this entire factor-by-factor strategy is impossible under the current
criterion.

\subsection{A valuation-lattice coefficient bound}

\begin{lemma}[Cramer bound for a rational multiplicative span]\label{lem:cramer}
Let $r_1,\ldots,r_m\in\Q_{>0}$, let
$w_i=\val(r_i)$, and suppose $e_2\in\Span_\Q\{w_1,\ldots,w_m\}$.  If every coordinate of
every $w_i$ has absolute value at most $E\ge1$, then there exist
$t\le m$, indices $i_1,\ldots,i_t$, a nonzero integer $d$, and integers
$c_1,\ldots,c_t$ such that
\begin{equation}\label{eq:cramer-rel}
  d e_2=\sum_{j=1}^t c_jw_{i_j}
\end{equation}
and
\begin{equation}\label{eq:cramer-bound}
  |c_j|\le (t-1)^{(t-1)/2}E^{t-1},
\end{equation}
with the right-hand side interpreted as $1$ for $t=1$.
\end{lemma}

\begin{proof}
Choose $t$ independent columns spanning $e_2$.  Some $t\times t$ prime-row minor $B$ is
invertible.  Such a minor necessarily contains the row indexed by $2$: otherwise the
restricted right-hand side is zero, so the uniquely determined coefficient vector is zero,
contrary to $e_2\ne0$.  Cramer's rule gives \eqref{eq:cramer-rel} with
$d=\det B$ and the $c_j$ equal to cofactors in the $2$-row column.  Hadamard's inequality on
each $(t-1)\times(t-1)$ minor gives \eqref{eq:cramer-bound}.
\end{proof}

\subsection[No rational tuple can span log 2]{No rational tuple can span $\log 2$}

\begin{theorem}[Rational-transfer impossibility]\label{thm:rational-transfer}
Let $z_1,\ldots,z_m\in\Q\cap(0,1)$ be pairwise distinct.  Suppose that they admit a
specialization of Kawashima's weight-two criterion with positive defect at the real place.
Put
\[
  \ell_i=\Li_1(z_i)=-\log(1-z_i).
\]
Then
\begin{equation}\label{eq:log2-notspan}
  \log2\notin\Span_\Q\{\ell_1,\ldots,\ell_m\}.
\end{equation}
The same conclusion holds with $2$ replaced by any fixed rational number after changing only
the harmless explicit constant in the proof.
\end{theorem}

\begin{proof}
Set
\[
 a=-\log\max_i z_i,\qquad P=\frac1m\sum_i h(z_i),
 \qquad D=\D_m,\qquad C=\Kc_m.
\]
By \Cref{cor:efficiency},
\begin{equation}\label{eq:aPconditions}
 a>C,\qquad P<\left(1+\frac1D\right)a.
\end{equation}
Write $r_i=(1-z_i)^{-1}$.  If $z_i=u_i/v_i$ in lowest terms with
$0<u_i<v_i$, then $r_i=v_i/(v_i-u_i)$ is also in lowest terms and
\begin{equation}\label{eq:equal-heights}
 h(r_i)=\log v_i=h(z_i).
\end{equation}
Because all heights are nonnegative,
\[
 h(r_i)=h(z_i)\le mP<m\left(1+\frac1D\right)a.
\]
Every prime valuation of $r_i$ is consequently bounded by
\begin{equation}\label{eq:Ebound}
 |v_p(r_i)|\le\frac{h(r_i)}{\log2}<3ma=:E.
\end{equation}
Also $a>C>\log2$, so $z_i\le e^{-a}<1/2$ and
\begin{equation}\label{eq:ell-small}
 0<\ell_i=-\log(1-z_i)\le\frac{z_i}{1-z_i}\le2e^{-a}.
\end{equation}

Assume for contradiction that $\log2$ belongs to the rational span of the $\ell_i$.
Exponentiating after clearing denominators shows that $e_2$ belongs to the rational span of
the valuation vectors $w_i=\val(r_i)$.  Apply \Cref{lem:cramer}.  Taking logarithms in
\eqref{eq:cramer-rel} and using \eqref{eq:ell-small} gives
\begin{align}
 \log2
 &\le |d|\log2
  =\left|\sum_{j=1}^t c_j\ell_{i_j}\right| \\
 &\le 2t(t-1)^{(t-1)/2}(3ma)^{t-1}e^{-a} \\
 &\le 2m\,m^{(m-1)/2}(3ma)^{m-1}e^{-a}.
 \label{eq:master-contradiction}
\end{align}
It remains only to check that the last expression is smaller than $\log2$ whenever
$a>C$.

The function $(3ma)^{m-1}e^{-a}$ is decreasing for $a>m-1$, and $C>2m^2>m-1$.
Furthermore
\begin{equation}\label{eq:Cupperlower}
 2m^2+4m< C<4(m+1)^2.
\end{equation}
The upper bound follows by writing $y=m+1\ge2$ and observing that
$(2-\log2)y^2+\log2-3\log y>0$; the lower bound is immediate from
$C=(2+\log2)(m^2+2m)+3\log(m+1)+2$.
At $a=C$, the logarithm of the right-hand side of
\eqref{eq:master-contradiction}, divided by $\log2$, is at most
\[
 \log(3m)+(m-1)\bigl(\log48+3.5\log m\bigr)-2m^2-4m.
\]
For $m=1$ the inequality is immediate.  For real $m\ge2$, the negative of this expression is
positive at $m=2$ and has positive derivative; its derivative itself is increasing because
$4-3.5/m-2.5/m^2>0$.  Hence the expression is negative for every integer $m\ge2$.
This contradicts \eqref{eq:master-contradiction}.
\end{proof}

\begin{provedbox}
The theorem is universal in $m$.  It does not merely say that the obvious representations of
$\log2$ fail; it says that \emph{every rational tuple satisfying Kawashima's own height
hypothesis} is too arithmetically close to $1$ for its valuation lattice to recover the fixed
prime direction $e_2$.
\end{provedbox}

\subsection{The full formal quadratic transfer also fails}

Let $W=\Span_\Q\{\val((1-z_i)^{-1}):1\le i\le m\}$.  Any factor-by-factor rational
logarithm identity places the resulting quadratic forms in $\operatorname{Sym}^2(W)$.
The tensor-rank audit from \Cref{lem:tensor-rank} now gives a stronger corollary.

\begin{corollary}[No formal rational quadratic transfer]\label{cor:quadratic-transfer}
Assume a nonintegral Alaoglu--Erd\H{o}s counterexample and let $T$ be its prime-valuation
tensor \eqref{eq:tensor}.  There is no rational tuple $z_1,\ldots,z_m\in(0,1)$ such that
\begin{enumerate}[label=\textup{(\roman*)}]
\item a specialization of Kawashima's weight-two criterion at these points has positive defect;
\item $T\in\operatorname{Sym}^2(W)$ for
$W=\Span_\Q\{\val((1-z_i)^{-1})\}$.
\end{enumerate}
In particular, one cannot settle \eqref{eq:wall} by expressing its four logarithms through
rational $\Li_1$ identities and then invoking Kawashima's product independence.
\end{corollary}

\begin{proof}
Every tensor in $\operatorname{Sym}^2(W)$ has contraction range contained in $W$.
By \Cref{lem:tensor-rank}, the range of $T$ is
$S=\Span\{u,v,e_2,e_3\}$, so $e_2\in W$.  Applying $\lambda$ gives
$\log2$ in the rational span of the numbers $\log(1-z_i)^{-1}=\Li_1(z_i)$, contradicting
\Cref{thm:rational-transfer}.
\end{proof}

The corollary leaves a narrow but genuine escape hatch: one could use algebraic points,
unbounded number-field degree, and cross-embedding or norm identities that do not reduce to a
rational prime-valuation tensor inside one fixed field.  Sections below show why the obvious
versions of that escape hatch also encounter strong barriers.

\section{Fixed multiplicative groups: effective exceptional finiteness}

The projective defect bound combines cleanly with the effective approximation theorem of
Calegari--Dimitrov--Tang.

\begin{theorem}[Fixed-group Pad\'e tuples are effectively finite]\label{thm:fixed-group}
Fix a number field $K$, a place $v$, a finitely generated subgroup
$\Gamma<K^\times$, and an integer $m\ge1$.  Consider pairwise distinct
$\gamma_1,\ldots,\gamma_m\in\Gamma$ such that
\[
  z_i=1-\gamma_i,\qquad |z_i|_v<1.
\]
There are only finitely many such tuples for which the points $z_i$ admit a specialization of
Kawashima's weight-two criterion with positive defect.  Their heights are effectively bounded
in terms of $(K,v,\Gamma,m)$.
\end{theorem}

\begin{proof}
Let $a=a_v(\mathbf z)$ and $P=\overline h(\mathbf z)$.  Positivity and
\eqref{eq:Pupper} imply, for every $i$,
\[
 h(z_i)\le mP<\frac{m(\D_m+1)}{\D_m}a.
\]
Hence, with
\[
 \varepsilon_m=\frac{\D_m}{2m(\D_m+1)},
\]
we have $a>2\varepsilon_m h(z_i)$.  The standard height inequalities give
$|h(\gamma_i)-h(z_i)|\le\log2$.  If $h(\gamma_i)\le2\log2$, it is already bounded.
Otherwise
\[
 -\log|1-\gamma_i|_v=-\log|z_i|_v\ge a
   >\varepsilon_m h(\gamma_i),
\]
so
\[
 |1-\gamma_i|_v<H(\gamma_i)^{-\varepsilon_m}.
\]
The effective approximation theorem of \cite{CDTHolonomy}, applied with $A=1$, bounds
$h(\gamma_i)$ effectively.  Northcott finiteness in the fixed field $K$ then leaves only
finitely many tuples.
\end{proof}

\begin{remark}
The theorem does not prove Alaoglu--Erd\H{o}s.  Under a hypothetical counterexample the group
of $S$-units generated by the finitely many primes dividing $6MA$ is fixed, but the effective
exceptional set may be nonempty and its defining set $S$ is itself unknown.  What the theorem
does prove is that no asymptotic sequence inside a fixed $S$-unit group can drive Kawashima's
defect positive.  Any such proof would have to come from a finite exceptional configuration or
from fields/groups whose arithmetic complexity changes with the construction.
\end{remark}

\section{The near-unit unimodular normal form}

\subsection{Exact construction}

Let
\[
  \rho=\frac{\log3}{\log2}
\]
and let $p_n/q_n$ be its continued-fraction convergents.  Define
\[
  u_n=2^{p_n}3^{-q_n},\qquad U_n=M^{p_n}A^{-q_n}=u_n^x.
\]
Let $\varepsilon_n\in\{\pm1\}$ orient the first unit above $1$:
\[
  r_n=u_n^{\varepsilon_n}>1,\qquad R_n=U_n^{\varepsilon_n}=r_n^x>1.
\]
Consecutive convergents satisfy
$p_nq_{n+1}-p_{n+1}q_n=\pm1$.  Therefore $(r_n,r_{n+1})$ is a multiplicative basis of
$\langle2,3\rangle$ and $(R_n,R_{n+1})$ is a multiplicative basis of
$\langle M,A\rangle$.  Moreover
\begin{equation}\label{eq:microdet}
 \det\begin{pmatrix}
  \log r_n&\log r_{n+1}\\
  \log R_n&\log R_{n+1}
 \end{pmatrix}=0
\end{equation}
and all four units tend to $1$.

The usual convergent inequalities give
\begin{equation}\label{eq:conv-bounds}
 \frac{\log2}{q_n+q_{n+1}}
 <|\log u_n|<\frac{\log2}{q_{n+1}}.
\end{equation}
Thus a counterexample is equivalent to infinitely many exact weight-two determinant zeros in
an arbitrarily small neighborhood of $(1,1,1,1)$.

\subsection{Why mere proximity is not enough}

For a positive rational unit $r>1$, put
\[
 z(r)=1-r^{-1},\qquad
 \eta(r)=\frac{-\log z(r)}{h(z(r))}.
\]
If $r=N/D$ is reduced, then $z(r)=(N-D)/N$ is reduced and
$h(z(r))=\log N=h(r)$.  The ratios $(m+1)/m$ appearing in the product theorem of
Calegari--Dimitrov--Tang have numerator gap $1$ and therefore
\[
 \eta\!\left(\frac{m+1}{m}\right)=1.
\]
They lie on the maximal-efficiency boundary.

The convergent units are entirely different.  A lower bound for a nonzero linear form in
$\log2$ and $\log3$ gives effective constants $c,C>0$ such that
\[
 |p\log2-q\log3|>c\,\max\{|p|,|q|\}^{-C}.
\]
Together with \eqref{eq:conv-bounds}, this implies
\[
 -\log z(r_n)=O(\log q_n),\qquad h(z(r_n))\asymp q_n,
\]
so
\begin{equation}\label{eq:eta-to-zero}
  \eta(r_n)\longrightarrow0.
\end{equation}
The same conclusion holds for $R_n$.  Indeed $\log R_n=x\log r_n$, while the valuation
vectors $p_n\val(M)-q_n\val(A)$ have norm $\asymp q_n$ because $M$ and $A$ are
multiplicatively independent.

\begin{provedbox}
The natural microlocal form of a counterexample lives in the \emph{low-efficiency} regime
$\eta\to0$.  Kawashima's present criterion requires average efficiency
$>1-1/(m+1)^2$, and the proved arithmetic-holonomy product theorem treats the unit-gap locus
$\eta=1$.  The missing theorem must therefore cross an arithmetic-efficiency gap, not merely a
geometric ``points close to one'' gap.
\end{provedbox}

\subsection{A numerical sample}

\begin{center}
\begin{tabular}{@{}rrrrr@{}}
\toprule
$n$ & $p_n$ & $q_n$ & $|p_n\log2-q_n\log3|$ & height scale \\
\midrule
0 & 1 & 1 & $4.0547\times10^{-1}$ & $1$ \\
1 & 2 & 1 & $2.8768\times10^{-1}$ & $2$ \\
2 & 3 & 2 & $1.1778\times10^{-1}$ & $3$ \\
3 & 8 & 5 & $5.2116\times10^{-2}$ & $8$ \\
4 & 19 & 12 & $1.3551\times10^{-2}$ & $19$ \\
5 & 65 & 41 & $1.1463\times10^{-2}$ & $65$ \\
6 & 84 & 53 & $2.0881\times10^{-3}$ & $84$ \\
7 & 485 & 306 & $1.0222\times10^{-3}$ & $485$ \\
8 & 1054 & 665 & $4.3654\times10^{-5}$ & $1054$ \\
9 & 24727 & 15601 & $1.8195\times10^{-5}$ & $24727$ \\
\bottomrule
\end{tabular}
\end{center}

The relevant decay is polynomial in the exponent size, while high-efficiency Pad\'e criteria
need a gap essentially as small as the reciprocal arithmetic denominator.

\section{Why the remaining nearby tools do not transfer}

\subsection[p-adic specialization]{$p$-adic specialization}

\Cref{thm:projective-defect} applies just as well at a finite place, and in some examples a
$p$-adic local contribution can be much more favorable.  However Kawashima's conclusion at a
finite place is linear independence of $p$-adic polylogarithm values in $K_v$.  The real
relation \eqref{eq:wall} does not automatically become a relation between $p$-adic logarithms.
A comparison principle strong enough to transfer that relation would itself be a new period
or motivic theorem at approximately the same transcendence level as the original problem.
Thus choosing a finite place changes the values being proved independent; it does not simply
improve the height estimate for the real values.

\subsection{Powers of one logarithm}

Kawashima's appendix constructs new Pad\'e-type approximants for
\[
 \log(1-1/z),\ldots,\log^m(1-1/z)
\]
and proves the relevant functional determinant is nonzero.  This is valuable structural
input, but \eqref{eq:wall} is a cross-product of logarithms at different algebraic arguments.
Moreover the paper explicitly leaves the Poincar\'e-type recurrence asymptotics needed for
sharp numerical independence measures to future work.  A one-point power theorem cannot
collapse the multiplicative rank of $2$ and $3$ to one.

\subsection{The proved holonomy product theorem}

The theorem of Calegari--Dimitrov--Tang proves a genuine mixed product result, but specifically
for
\[
 \log\!\left(1+\frac1m\right)\log\!\left(1+\frac1n\right)
\]
with $m/n$ close to $1$.  A convergent unit $2^p/3^q$ generally has numerator gap
$|2^p-3^q|\gg1$, so it is not of the form $(m+1)/m$.  Pairwise irrationality of individual
products also does not imply joint linear independence of all products arising after a long
telescoping decomposition.  The theorem reaches weight two, but on a high-efficiency locus
disjoint from the canonical counterexample normal form.

\section{What a successful Pad\'e continuation would have to change}

The no-go theorems above identify three requirements that are absent from the current
criterion.

\subsection{Direct determinant approximation rather than monomial independence}

Kawashima proves the independence of a very large family: all polylogarithm monomials of total
weight at most two.  For $m$ points this family has size $(m+1)^2$, and the constant
$\Kc_m\asymp(2+\log2)m^2$ records that dimensional cost.  Alaoglu--Erd\H{o}s needs only one
specific determinant.  A tailored construction should approximate
\[
 \Delta=\log M\log3-\log A\log2
\]
directly instead of paying for every monomial.

\subsection{Separate local denominators}

The common parameter $\beta$ forces all points through one affine denominator and charges the
heights of both the coordinates and their vector.  A plausible replacement is a genuinely
mixed-denominator Hermite--Pad\'e system with one local denominator for each point, followed by
adelic cancellation \emph{before} a common denominator is cleared.  Merely rewriting the same
points with a larger common denominator is impossible by \eqref{eq:scaling-law}.

\subsection{Valuation-lattice cancellation}

The rational-transfer theorem shows that individual high-efficiency logarithms are too small
to recover a fixed prime direction with the polynomially bounded coefficients supplied by a
valuation lattice.  A successful construction must therefore exploit cancellation at the
level of the \emph{four-log determinant}: prime factors appearing in one Pad\'e row must cancel
against prime factors in another row.  This is precisely the type of cross-row gain that the
ordinary one-row height defect discards.

The following statement isolates a sufficient target without pretending that it is already
proved.

\begin{conjecture}[Mixed-denominator determinant criterion]\label{conj:mixed}
For positive rationals $r_1,r_2,s_1,s_2$ such that each pair is multiplicatively independent,
there is a Pad\'e construction for
\[
 \log r_1\log s_2-\log r_2\log s_1
\]
whose denominator estimate is governed by the determinantal divisors of the two prime-valuation
matrices, including their saturation indices, rather than by the sum of the four individual
affine heights.  The construction proves that the displayed determinant is nonzero unless a
rational row/column change produces a zero logarithmic coefficient.
\end{conjecture}

\begin{proposition}\label{prop:mixed-implies}
\Cref{conj:mixed} implies the Alaoglu--Erd\H{o}s conjecture.
\end{proposition}

\begin{proof}
Assume a counterexample.  Use consecutive convergents as in \Cref{eq:microdet}.  The pairs
$(r_n,r_{n+1})$ and $(R_n,R_{n+1})$ are unimodular changes of the respective rank-two
valuation-lattice generators, while
their logarithmic determinant is zero.  The exceptional alternative in
\Cref{conj:mixed} would give a zero logarithmic coefficient after rational basis changes.  By
injectivity of the rational logarithm map, that coefficient is zero exactly when one of the
corresponding multiplicative basis elements is $1$, impossible.  Hence the determinant must be
nonzero, contradicting \eqref{eq:microdet}.
\end{proof}

The conjectural criterion is deliberately narrower than Four Exponentials in arbitrary
algebraic settings: it asks only for positive rationals and rank-two prime-valuation
lattices.  It is also deliberately stronger than the existing Pad\'e theorems in the one place
where the present problem needs strength: low-efficiency, cross-row denominator cancellation.

\section{\texorpdfstring{Directed-rounding extension through $2^{29}$}{Directed-rounding extension through 2 power 29}}\label{sec:scan}

\subsection{Verifier contract}

Let
\[
  \rho=\frac{\log3}{\log2}.
\]
For each positive integer $m$ that is not a power of two, the verifier seeks an integer $n$
with
\[
  n<m^\rho<n+1.
\]
It never evaluates $m^\rho$.  At $192$-bit MPFR precision it computes directed enclosures of
the needed logarithms and proves the two strict inequalities through positive lower bounds for
\begin{align}
 \delta_-&=\log(m)\log3-\log(n)\log2,\\
 \delta_+&=\log(n+1)\log2-\log(m)\log3.
\end{align}
A long-double calculation proposes $n$, but the proposal has no logical role: acceptance
requires both MPFR lower bounds to be strictly positive.  A locator error can therefore create
a reported failure, not a false certificate.

The exact C source was extracted from the supplied report without modification.  Its SHA-256
digest is
\begin{center}
\sha{f888c0d89be9366f01e2d26f339dcbe446556d39d4c5dbcf8f579ec0d0052e1d}.
\end{center}
This identity matches the source digest already printed in the supplied report.

\subsection{New range and aggregate result}

The new audit covers two consecutive bit ranges,
\[
 [2^{27},2^{29})=[134{,}217{,}728,536{,}870{,}912),
\]
partitioned into $96$ chunks of width $2^{22}$.  The powers $2^{27}$ and $2^{28}$ are skipped.
Consequently the exact number of tested inputs is
\[
 (2^{29}-2^{27})-2=402{,}653{,}182.
\]
The first $32$ chunks cover $[2^{27},2^{28})$; the next $64$ cover the entire range
$[2^{28},2^{29})$.

\begin{center}
\begin{tabular}{@{}lrrr@{}}
\toprule
Range & tested nonpowers & failures & summed chunk time (s) \\
\midrule
$[2^{27},2^{28})$ & $134{,}217{,}727$ & $0$ & $395.843$ \\
$[2^{28},2^{29})$ & $268{,}435{,}455$ & $0$ & $792.962$ \\
\midrule
Combined new audit & $402{,}653{,}182$ & $0$ & $1{,}188.805$ \\
\bottomrule
\end{tabular}
\end{center}

Across both new ranges, the smallest certified lower-side margin was
\begin{equation}\label{eq:scan-lower}
 \delta_->4.9141462461940518\times10^{-24}
 \quad\text{at}\quad
 (m,n)=(388{,}909{,}543,41{,}170{,}668{,}399{,}174),
\end{equation}
and the smallest certified upper-side margin was
\begin{equation}\label{eq:scan-upper}
 \delta_+>5.2103670350428438\times10^{-23}
 \quad\text{at}\quad
 (m,n)=(501{,}373{,}708,61{,}578{,}600{,}753{,}798).
\end{equation}
Both extremal records were replayed with one thread at $1024$-bit MPFR precision.  The
corresponding directed lower bounds were unchanged to the displayed $17$ digits.  As a
separate implementation check, a $320$-decimal-digit evaluation gave
\[
 \delta_-=4.914146246194051848982433252362132380596\ldots\times10^{-24}
\]
for \eqref{eq:scan-lower}, and
\[
 \delta_+=5.210367035042843821555189929112177636068\ldots\times10^{-23}
\]
for \eqref{eq:scan-upper}.  This decimal replay is corroborative only; the certificate is the
positive directed MPFR lower bound.

The executable used GCC~14.2.0, MPFR~4.2.2, $192$-bit precision, and five OpenMP threads for
each chunk.  The source and executable SHA-256 digests were, respectively,
\begin{align*}
 &\sha{f888c0d89be9366f01e2d26f339dcbe446556d39d4c5dbcf8f579ec0d0052e1d},\\
 &\sha{85f2acca22ec3ea4bf783a2e13356a1e11d9faff3b95a4d1c79a7cced69deceb}.
\end{align*}
The concatenated logs for the two blocks had SHA-256 digests
\begin{align*}
 [2^{27},2^{28})&:\quad
 \sha{3865fb3d99d3a6df4ab4ac101942781dea1065f38414462fb2e25b0ba4ef16ce},\\
 [2^{28},2^{29})&:\quad
 \sha{6899851a9d78bca4e7527011c10089aaede12f1f93a9584efdcaa0d5d56cebe2}.
\end{align*}
The two $1024$-bit replay logs had digests
\begin{align*}
 &\sha{333b015a3ac541f34a4c14197dc67d03d5211f7bc0828abd346695790a174555},\\
 &\sha{de0eac89dcc88c63f080e8714fa5e93f0485beb8e913e87ab705316b80be619b},
\end{align*}
and the independent $320$-digit replay file had digest
\[
 \sha{791d592dc7140674e8300c2ae65db6ef6be9f539b1f92a05a8b6f5f4c37be223}.
\]

\begin{compbox}
The result is a software-level directed-rounding certificate, not a proof-assistant kernel
certificate.  It depends on the stated MPFR implementation, compiler, ABI declarations, and
executed binary.  Those trust assumptions are materially stronger than ordinary floating-point
sampling but materially weaker than a formal proof.
\end{compbox}

\subsection{Consequence for the exponent}

\begin{proposition}[Finite exclusion through $29$]\label{prop:x29}
At the trust level of the directed-rounding computation, any nonintegral counterexample
satisfies $x>29$.
\end{proposition}

\begin{proof}
Suppose $x<29$ and $M=2^x\in\N$.  Then $M<2^{29}$.  If $M$ is a power of two, injectivity of
the real exponential gives $x\in\Z$.  Otherwise the inherited scan below $2^{27}$ together
with the two new ranges proves that $M^\rho=3^x$ lies strictly between consecutive integers,
contradicting the hypothesis that $3^x$ is an integer.  The endpoint $x=29$ is integral.
\end{proof}

\section{Synthesis}

The post-report Pad\'e literature gets closer to the correct transcendence weight than any of
the classical one-logarithm tools, but the new analysis identifies an exact mismatch.

\begin{enumerate}[label=\textbf{\alph*.}]
\item A counterexample is naturally converted by continued fractions into determinant zeros
      among rational units arbitrarily close to $1$.
\item Those units have arithmetic efficiency tending to $0$: their gaps are only polynomially
      small in the exponent size, while their rational heights are exponential in that size.
\item Kawashima's current theorem requires average efficiency
      $>1-1/(m+1)^2$ and a local depth exceeding a quadratic dimension constant.
\item The proven arithmetic-holonomy product theorem occupies the extreme unit-gap locus of
      efficiency $1$.
\item Rational algebraic transfer cannot bridge the gap: the valuation lattice bounds the
      required coefficients polynomially, while all qualifying logarithms are exponentially
      small.  This is the content of \Cref{thm:rational-transfer}.
\end{enumerate}

The remaining route is therefore not ``find a cleverer common denominator.''  It is to build a
weight-two approximant that sees the prime-valuation determinant of the \emph{whole} expression
and cancels denominators across rows before individual heights are charged.  That target is
technically concrete, but it is not presently a theorem.

\section{Final status}

\begin{conditionalbox}
No unconditional proof of
\[
  2^x,3^x\in\Z\quad\Longrightarrow\quad x\in\Z
\]
is contained in this report.  The exact unresolved statement remains the nonvanishing of the
weight-two logarithmic determinant \eqref{eq:wall}, equivalently the corresponding restricted
matrix-coefficient assertion.
\end{conditionalbox}

The rigorous advances are nonetheless substantive:

\begin{itemize}
\item a new projective upper bound for the 2026 polylogarithm-product criterion;
\item a quantitative arithmetic-efficiency threshold with sharp dependence on the number of
      points;
\item complete no-go results for direct points, radical splitting, ordinary telescoping, and
      common scaling;
\item a universal rational-transfer impossibility theorem, plus a tensor-level corollary for
      the formal Alaoglu--Erd\H{o}s quadratic relation;
\item effective finiteness of every fixed-group Pad\'e attempt;
\item a precise diagnosis of the high-efficiency/low-efficiency gap between current theorems
      and the canonical counterexample normal form;
\item an independently rerun directed-rounding exclusion through $2^{29}$.
\end{itemize}

These results eliminate a broad family of plausible but ultimately nonviable proof strategies
and leave a narrower, technically explicit research target: direct determinant Pad\'e
approximation with cross-row valuation cancellation.

\appendix

\section{Elementary inequalities used in the rational-transfer theorem}

For completeness, this appendix records the two estimates compressed in
\Cref{thm:rational-transfer}.

Let $D=m^2+2m$ and
\[
 C=(2+\log2)D+3\log(m+1)+2.
\]
Then $C>2m^2+4m$ is immediate.  For the upper bound, put $y=m+1\ge2$:
\begin{align*}
 4(m+1)^2-C
 &=4y^2-(2+\log2)(y^2-1)-3\log y-2\\
 &=(2-\log2)y^2+\log2-3\log y.
\end{align*}
The last expression is positive at $y=2$, and its derivative
$2(2-\log2)y-3/y$ is positive for $y\ge2$.  Hence
\[
 2m^2+4m<C<4(m+1)^2.
\]

For $m\ge2$, define
\[
 G(m)=2m^2+4m-\log(3m)-(m-1)(\log48+3.5\log m).
\]
Directly, $G(2)>7.9$.  Moreover
\[
 G'(m)=4m+\frac12+\frac{2.5}{m}-\log48-3.5\log m>0
\]
for $m\ge2$, because its derivative
$4-3.5/m-2.5/m^2$ is already positive at $m=2$.  Thus $G(m)>0$ for all integers $m\ge2$.
This is the last numerical inequality needed to make
\eqref{eq:master-contradiction} impossible.

\section{Scan chunk manifest}

The table records every chunk.  Each interval has width $2^{22}$; ``start'' is its
inclusive lower endpoint.  A record $d\;(m)$ means that the chunk minimum was the directed
lower bound $d$ attained at that value of $m$.  The corresponding integer $n$ is retained in
the hashed raw log; the two global $n$ values are displayed in
\eqref{eq:scan-lower}--\eqref{eq:scan-upper}.

\begingroup
\scriptsize
\setlength{\tabcolsep}{3.2pt}
\begin{longtable}{@{}c r r >{\raggedleft\arraybackslash}p{34mm} >{\raggedleft\arraybackslash}p{34mm} r@{}}
\caption{Directed-rounding chunk manifest for $[2^{27},2^{29})$.}\label{tab:scan-manifest}\\
\toprule
ID & start & tested & minimum $\delta_-$ $(m)$ & minimum $\delta_+$ $(m)$ & seconds \\
\midrule
\endfirsthead
\multicolumn{6}{c}{\tablename\ \thetable\ continued}\\
\toprule
ID & start & tested & minimum $\delta_-$ $(m)$ & minimum $\delta_+$ $(m)$ & seconds \\
\midrule
\endhead
\midrule
\multicolumn{6}{r}{continued on next page}\\
\endfoot
\bottomrule
\endlastfoot
27.00 & 134{,}217{,}728 & 4{,}194{,}303 & $4.284e-20$ $(136{,}640{,}874)$ & $3.389e-20$ $(135{,}357{,}389)$ & 11.825 \\
27.01 & 138{,}412{,}032 & 4{,}194{,}304 & $2.815e-20$ $(138{,}540{,}675)$ & $2.847e-21$ $(142{,}153{,}616)$ & 12.013 \\
27.02 & 142{,}606{,}336 & 4{,}194{,}304 & $1.061e-20$ $(142{,}695{,}401)$ & $6.273e-22$ $(146{,}333{,}091)$ & 12.881 \\
27.03 & 146{,}800{,}640 & 4{,}194{,}304 & $3.593e-21$ $(148{,}583{,}161)$ & $1.143e-20$ $(148{,}016{,}970)$ & 11.572 \\
27.04 & 150{,}994{,}944 & 4{,}194{,}304 & $3.782e-21$ $(153{,}415{,}953)$ & $4.828e-20$ $(154{,}129{,}681)$ & 13.257 \\
27.05 & 155{,}189{,}248 & 4{,}194{,}304 & $1.677e-20$ $(159{,}254{,}759)$ & $1.606e-20$ $(158{,}382{,}018)$ & 12.778 \\
27.06 & 159{,}383{,}552 & 4{,}194{,}304 & $3.469e-20$ $(162{,}200{,}015)$ & $2.021e-20$ $(160{,}095{,}403)$ & 11.821 \\
27.07 & 163{,}577{,}856 & 4{,}194{,}304 & $3.847e-21$ $(166{,}033{,}450)$ & $7.569e-21$ $(166{,}969{,}115)$ & 12.110 \\
27.08 & 167{,}772{,}160 & 4{,}194{,}304 & $8.665e-21$ $(171{,}251{,}435)$ & $1.465e-20$ $(170{,}443{,}795)$ & 11.822 \\
27.09 & 171{,}966{,}464 & 4{,}194{,}304 & $4.262e-20$ $(173{,}816{,}254)$ & $1.433e-20$ $(175{,}647{,}581)$ & 11.661 \\
27.10 & 176{,}160{,}768 & 4{,}194{,}304 & $1.876e-20$ $(178{,}598{,}559)$ & $4.477e-21$ $(179{,}458{,}795)$ & 13.244 \\
27.11 & 180{,}355{,}072 & 4{,}194{,}304 & $5.679e-20$ $(180{,}765{,}336)$ & $5.445e-22$ $(183{,}569{,}533)$ & 12.511 \\
27.12 & 184{,}549{,}376 & 4{,}194{,}304 & $3.085e-22$ $(187{,}004{,}237)$ & $7.872e-21$ $(188{,}301{,}745)$ & 12.227 \\
27.13 & 188{,}743{,}680 & 4{,}194{,}304 & $1.102e-20$ $(192{,}461{,}540)$ & $1.798e-20$ $(191{,}799{,}137)$ & 12.614 \\
27.14 & 192{,}937{,}984 & 4{,}194{,}304 & $1.638e-21$ $(195{,}368{,}525)$ & $2.173e-20$ $(194{,}355{,}394)$ & 11.846 \\
27.15 & 197{,}132{,}288 & 4{,}194{,}304 & $8.763e-21$ $(197{,}987{,}777)$ & $1.177e-20$ $(198{,}330{,}250)$ & 12.451 \\
27.16 & 201{,}326{,}592 & 4{,}194{,}304 & $9.306e-22$ $(202{,}612{,}057)$ & $9.219e-21$ $(203{,}389{,}040)$ & 13.492 \\
27.17 & 205{,}520{,}896 & 4{,}194{,}304 & $2.745e-21$ $(206{,}738{,}514)$ & $1.370e-21$ $(206{,}597{,}929)$ & 11.626 \\
27.18 & 209{,}715{,}200 & 4{,}194{,}304 & $1.383e-20$ $(213{,}551{,}311)$ & $8.366e-21$ $(213{,}290{,}074)$ & 13.091 \\
27.19 & 213{,}909{,}504 & 4{,}194{,}304 & $4.101e-22$ $(214{,}684{,}191)$ & $1.679e-21$ $(217{,}426{,}809)$ & 12.462 \\
27.20 & 218{,}103{,}808 & 4{,}194{,}304 & $3.638e-21$ $(218{,}691{,}392)$ & $2.873e-21$ $(218{,}721{,}921)$ & 12.370 \\
27.21 & 222{,}298{,}112 & 4{,}194{,}304 & $1.825e-20$ $(224{,}977{,}296)$ & $5.053e-21$ $(226{,}024{,}581)$ & 12.121 \\
27.22 & 226{,}492{,}416 & 4{,}194{,}304 & $1.646e-20$ $(226{,}553{,}247)$ & $2.228e-20$ $(230{,}684{,}216)$ & 13.078 \\
27.23 & 230{,}686{,}720 & 4{,}194{,}304 & $1.953e-22$ $(233{,}957{,}087)$ & $6.347e-21$ $(231{,}792{,}276)$ & 12.023 \\
27.24 & 234{,}881{,}024 & 4{,}194{,}304 & $3.812e-21$ $(235{,}808{,}239)$ & $3.750e-21$ $(236{,}568{,}448)$ & 12.491 \\
27.25 & 239{,}075{,}328 & 4{,}194{,}304 & $6.046e-21$ $(239{,}456{,}283)$ & $9.724e-21$ $(239{,}580{,}935)$ & 13.390 \\
27.26 & 243{,}269{,}632 & 4{,}194{,}304 & $3.971e-21$ $(244{,}524{,}864)$ & $5.006e-21$ $(245{,}994{,}981)$ & 11.525 \\
27.27 & 247{,}463{,}936 & 4{,}194{,}304 & $1.077e-20$ $(249{,}872{,}818)$ & $8.052e-21$ $(249{,}272{,}020)$ & 12.868 \\
27.28 & 251{,}658{,}240 & 4{,}194{,}304 & $1.929e-20$ $(251{,}728{,}391)$ & $4.421e-21$ $(254{,}182{,}804)$ & 12.215 \\
27.29 & 255{,}852{,}544 & 4{,}194{,}304 & $6.990e-21$ $(256{,}676{,}151)$ & $8.307e-22$ $(259{,}598{,}690)$ & 11.862 \\
27.30 & 260{,}046{,}848 & 4{,}194{,}304 & $4.051e-21$ $(263{,}613{,}699)$ & $9.409e-22$ $(261{,}264{,}827)$ & 12.792 \\
27.31 & 264{,}241{,}152 & 4{,}194{,}304 & $1.374e-21$ $(266{,}715{,}185)$ & $4.150e-21$ $(264{,}885{,}749)$ & 11.805 \\
28.00 & 268{,}435{,}456 & 4{,}194{,}303 & $1.438e-20$ $(270{,}769{,}256)$ & $1.276e-20$ $(269{,}492{,}939)$ & 12.071 \\
28.01 & 272{,}629{,}760 & 4{,}194{,}304 & $1.294e-20$ $(274{,}601{,}332)$ & $1.329e-20$ $(276{,}729{,}727)$ & 11.772 \\
28.02 & 276{,}824{,}064 & 4{,}194{,}304 & $1.336e-21$ $(278{,}937{,}929)$ & $1.473e-20$ $(279{,}209{,}135)$ & 13.839 \\
28.03 & 281{,}018{,}368 & 4{,}194{,}304 & $6.220e-21$ $(283{,}347{,}824)$ & $2.847e-21$ $(284{,}307{,}232)$ & 12.520 \\
28.04 & 285{,}212{,}672 & 4{,}194{,}304 & $3.573e-21$ $(285{,}482{,}152)$ & $1.521e-20$ $(288{,}577{,}614)$ & 12.190 \\
28.05 & 289{,}406{,}976 & 4{,}194{,}304 & $8.508e-21$ $(290{,}399{,}029)$ & $6.273e-22$ $(292{,}666{,}182)$ & 12.526 \\
28.06 & 293{,}601{,}280 & 4{,}194{,}304 & $3.593e-21$ $(297{,}166{,}322)$ & $3.235e-21$ $(297{,}393{,}168)$ & 12.063 \\
28.07 & 297{,}795{,}584 & 4{,}194{,}304 & $1.775e-21$ $(301{,}090{,}986)$ & $3.276e-20$ $(298{,}685{,}330)$ & 12.413 \\
28.08 & 301{,}989{,}888 & 4{,}194{,}304 & $1.725e-21$ $(304{,}348{,}427)$ & $3.239e-21$ $(303{,}439{,}436)$ & 13.260 \\
28.09 & 306{,}184{,}192 & 4{,}194{,}304 & $5.191e-22$ $(306{,}228{,}765)$ & $2.879e-20$ $(309{,}957{,}834)$ & 11.720 \\
28.10 & 310{,}378{,}496 & 4{,}194{,}304 & $3.887e-21$ $(313{,}261{,}071)$ & $5.204e-21$ $(313{,}114{,}243)$ & 12.893 \\
28.11 & 314{,}572{,}800 & 4{,}194{,}304 & $1.048e-20$ $(316{,}959{,}535)$ & $3.653e-22$ $(314{,}764{,}391)$ & 12.743 \\
28.12 & 318{,}767{,}104 & 4{,}194{,}304 & $1.666e-20$ $(321{,}040{,}811)$ & $2.032e-21$ $(320{,}803{,}384)$ & 12.558 \\
28.13 & 322{,}961{,}408 & 4{,}194{,}304 & $1.378e-20$ $(324{,}726{,}584)$ & $2.016e-20$ $(325{,}862{,}257)$ & 12.591 \\
28.14 & 327{,}155{,}712 & 4{,}194{,}304 & $9.296e-21$ $(330{,}041{,}528)$ & $1.951e-21$ $(327{,}923{,}577)$ & 12.732 \\
28.15 & 331{,}350{,}016 & 4{,}194{,}304 & $3.847e-21$ $(332{,}066{,}900)$ & $6.512e-21$ $(334{,}016{,}730)$ & 12.549 \\
28.16 & 335{,}544{,}320 & 4{,}194{,}304 & $1.113e-21$ $(338{,}945{,}756)$ & $2.743e-21$ $(336{,}494{,}503)$ & 12.405 \\
28.17 & 339{,}738{,}624 & 4{,}194{,}304 & $7.650e-21$ $(343{,}687{,}066)$ & $1.989e-21$ $(343{,}467{,}223)$ & 12.313 \\
28.18 & 343{,}932{,}928 & 4{,}194{,}304 & $2.000e-21$ $(347{,}267{,}147)$ & $2.772e-21$ $(346{,}761{,}393)$ & 12.206 \\
28.19 & 348{,}127{,}232 & 4{,}194{,}304 & $1.437e-22$ $(349{,}972{,}684)$ & $1.433e-20$ $(351{,}295{,}162)$ & 13.603 \\
28.20 & 352{,}321{,}536 & 4{,}194{,}304 & $5.143e-22$ $(353{,}661{,}837)$ & $2.833e-21$ $(353{,}990{,}351)$ & 11.765 \\
28.21 & 356{,}515{,}840 & 4{,}194{,}304 & $3.081e-21$ $(356{,}895{,}426)$ & $5.735e-22$ $(357{,}765{,}549)$ & 11.813 \\
28.22 & 360{,}710{,}144 & 4{,}194{,}304 & $8.141e-22$ $(364{,}737{,}335)$ & $1.316e-22$ $(362{,}062{,}452)$ & 13.552 \\
28.23 & 364{,}904{,}448 & 4{,}194{,}304 & $9.240e-22$ $(365{,}075{,}693)$ & $3.223e-22$ $(365{,}778{,}901)$ & 11.562 \\
28.24 & 369{,}098{,}752 & 4{,}194{,}304 & $7.157e-22$ $(371{,}893{,}339)$ & $8.105e-22$ $(371{,}246{,}341)$ & 11.810 \\
28.25 & 373{,}293{,}056 & 4{,}194{,}304 & $3.085e-22$ $(374{,}008{,}474)$ & $2.529e-21$ $(376{,}055{,}882)$ & 13.831 \\
28.26 & 377{,}487{,}360 & 4{,}194{,}304 & $8.880e-21$ $(378{,}129{,}221)$ & $6.429e-21$ $(381{,}433{,}624)$ & 12.044 \\
28.27 & 381{,}681{,}664 & 4{,}194{,}304 & $8.273e-22$ $(384{,}046{,}343)$ & $1.056e-20$ $(382{,}597{,}108)$ & 12.525 \\
28.28 & 385{,}875{,}968 & 4{,}194{,}304 & $4.914e-24$ $(388{,}909{,}543)$ & $8.857e-21$ $(386{,}125{,}278)$ & 12.275 \\
28.29 & 390{,}070{,}272 & 4{,}194{,}304 & $1.638e-21$ $(390{,}737{,}050)$ & $8.114e-22$ $(391{,}647{,}767)$ & 11.700 \\
28.30 & 394{,}264{,}576 & 4{,}194{,}304 & $2.526e-21$ $(395{,}220{,}414)$ & $5.912e-21$ $(398{,}186{,}911)$ & 12.558 \\
28.31 & 398{,}458{,}880 & 4{,}194{,}304 & $7.885e-24$ $(400{,}544{,}529)$ & $1.694e-20$ $(402{,}139{,}190)$ & 12.558 \\
28.32 & 402{,}653{,}184 & 4{,}194{,}304 & $9.306e-22$ $(405{,}224{,}114)$ & $4.711e-22$ $(405{,}024{,}519)$ & 12.028 \\
28.33 & 406{,}847{,}488 & 4{,}194{,}304 & $9.700e-22$ $(407{,}982{,}235)$ & $2.819e-21$ $(407{,}281{,}054)$ & 12.332 \\
28.34 & 411{,}041{,}792 & 4{,}194{,}304 & $2.745e-21$ $(413{,}477{,}028)$ & $1.293e-21$ $(412{,}276{,}830)$ & 11.854 \\
28.35 & 415{,}236{,}096 & 4{,}194{,}304 & $5.219e-21$ $(416{,}468{,}676)$ & $7.693e-22$ $(418{,}544{,}766)$ & 11.578 \\
28.36 & 419{,}430{,}400 & 4{,}194{,}304 & $1.954e-21$ $(422{,}398{,}112)$ & $3.436e-21$ $(423{,}295{,}624)$ & 12.982 \\
28.37 & 423{,}624{,}704 & 4{,}194{,}304 & $1.289e-21$ $(423{,}710{,}706)$ & $2.513e-22$ $(426{,}529{,}906)$ & 12.895 \\
28.38 & 427{,}819{,}008 & 4{,}194{,}304 & $4.101e-22$ $(429{,}368{,}382)$ & $4.883e-22$ $(429{,}282{,}395)$ & 13.442 \\
28.39 & 432{,}013{,}312 & 4{,}194{,}304 & $5.984e-21$ $(432{,}505{,}720)$ & $1.679e-21$ $(434{,}853{,}618)$ & 12.570 \\
28.40 & 436{,}207{,}616 & 4{,}194{,}304 & $1.685e-21$ $(436{,}314{,}750)$ & $2.873e-21$ $(437{,}443{,}842)$ & 11.926 \\
28.41 & 440{,}401{,}920 & 4{,}194{,}304 & $1.275e-21$ $(442{,}841{,}765)$ & $5.083e-21$ $(442{,}062{,}629)$ & 11.858 \\
28.42 & 444{,}596{,}224 & 4{,}194{,}304 & $4.070e-21$ $(447{,}235{,}367)$ & $7.158e-21$ $(447{,}720{,}068)$ & 13.405 \\
28.43 & 448{,}790{,}528 & 4{,}194{,}304 & $1.569e-21$ $(452{,}029{,}583)$ & $2.822e-22$ $(450{,}505{,}498)$ & 12.393 \\
28.44 & 452{,}984{,}832 & 4{,}194{,}304 & $4.410e-21$ $(455{,}437{,}571)$ & $1.014e-21$ $(454{,}240{,}771)$ & 11.531 \\
28.45 & 457{,}179{,}136 & 4{,}194{,}304 & $2.793e-21$ $(459{,}828{,}193)$ & $8.322e-21$ $(461{,}365{,}410)$ & 12.483 \\
28.46 & 461{,}373{,}440 & 4{,}194{,}304 & $1.747e-21$ $(463{,}276{,}633)$ & $6.347e-21$ $(463{,}584{,}552)$ & 11.568 \\
28.47 & 465{,}567{,}744 & 4{,}194{,}304 & $1.953e-22$ $(467{,}914{,}174)$ & $7.459e-22$ $(466{,}121{,}593)$ & 11.721 \\
28.48 & 469{,}762{,}048 & 4{,}194{,}304 & $3.067e-21$ $(471{,}111{,}503)$ & $3.750e-21$ $(473{,}136{,}896)$ & 13.749 \\
28.49 & 473{,}956{,}352 & 4{,}194{,}304 & $3.276e-21$ $(477{,}377{,}170)$ & $2.199e-22$ $(474{,}560{,}743)$ & 11.791 \\
28.50 & 478{,}150{,}656 & 4{,}194{,}304 & $5.934e-21$ $(480{,}943{,}527)$ & $6.086e-21$ $(478{,}837{,}280)$ & 11.811 \\
28.51 & 482{,}344{,}960 & 4{,}194{,}304 & $4.365e-22$ $(485{,}602{,}436)$ & $4.656e-22$ $(483{,}731{,}011)$ & 12.214 \\
28.52 & 486{,}539{,}264 & 4{,}194{,}304 & $5.845e-22$ $(488{,}578{,}708)$ & $2.469e-22$ $(490{,}595{,}136)$ & 12.000 \\
28.53 & 490{,}733{,}568 & 4{,}194{,}304 & $6.973e-21$ $(494{,}115{,}822)$ & $1.425e-21$ $(491{,}328{,}133)$ & 11.615 \\
28.54 & 494{,}927{,}872 & 4{,}194{,}304 & $4.455e-21$ $(495{,}724{,}767)$ & $1.792e-21$ $(496{,}559{,}778)$ & 14.097 \\
28.55 & 499{,}122{,}176 & 4{,}194{,}304 & $5.780e-22$ $(499{,}782{,}534)$ & $5.210e-23$ $(501{,}373{,}708)$ & 11.701 \\
28.56 & 503{,}316{,}480 & 4{,}194{,}304 & $4.533e-21$ $(504{,}996{,}574)$ & $9.370e-22$ $(504{,}517{,}808)$ & 11.720 \\
28.57 & 507{,}510{,}784 & 4{,}194{,}304 & $1.140e-21$ $(507{,}876{,}391)$ & $1.183e-21$ $(509{,}911{,}917)$ & 12.174 \\
28.58 & 511{,}705{,}088 & 4{,}194{,}304 & $6.990e-21$ $(513{,}352{,}302)$ & $4.274e-22$ $(514{,}216{,}147)$ & 12.008 \\
28.59 & 515{,}899{,}392 & 4{,}194{,}304 & $1.115e-21$ $(516{,}239{,}075)$ & $8.307e-22$ $(519{,}197{,}380)$ & 12.192 \\
28.60 & 520{,}093{,}696 & 4{,}194{,}304 & $6.023e-22$ $(521{,}145{,}596)$ & $9.409e-22$ $(522{,}529{,}654)$ & 14.640 \\
28.61 & 524{,}288{,}000 & 4{,}194{,}304 & $1.326e-21$ $(528{,}412{,}351)$ & $1.972e-21$ $(526{,}958{,}976)$ & 12.038 \\
28.62 & 528{,}482{,}304 & 4{,}194{,}304 & $4.360e-21$ $(528{,}890{,}012)$ & $9.307e-22$ $(530{,}509{,}954)$ & 11.560 \\
28.63 & 532{,}676{,}608 & 4{,}194{,}304 & $1.374e-21$ $(533{,}430{,}370)$ & $4.753e-21$ $(535{,}448{,}067)$ & 12.130 \\
\end{longtable}
\endgroup

\section{Reproducibility commands}

The source was compiled and the new range was scanned with commands equivalent to:

\begin{lstlisting}[style=shell]
gcc -O3 -fopenmp mpfr_scan.c /lib/x86_64-linux-gnu/libmpfr.so.6 \
    -lm -o mpfr_scan

for ((a=134217728; a<536870912; a+=4194304)); do
  b=$((a + 4194304))
  OMP_NUM_THREADS=5 ./mpfr_scan "$a" "$b" 192
done
\end{lstlisting}

The two global extremal records were rerun individually at $1024$ bits.  A high-precision
independent decimal evaluation of the true logarithmic margins is also reported in
\Cref{sec:scan}; it is not used in place of the directed MPFR inequalities.

\begin{thebibliography}{99}

\bibitem{MOAE}
MathOverflow community wiki,
\newblock \emph{If $2^x$ and $3^x$ are integers, must $x$ be as well?},
\newblock question 17560, accessed August 21, 2026.
\newblock \href{https://mathoverflow.net/questions/17560/if-2x-and-3x-are-integers-must-x-be-as-well}{MathOverflow 17560}.

\bibitem{AE1944}
L.~Alaoglu and P.~Erd\H{o}s,
\newblock \emph{On highly composite and similar numbers},
\newblock Transactions of the American Mathematical Society \textbf{56} (1944), 448--469.
\newblock \href{https://doi.org/10.1090/S0002-9947-1944-0011087-2}{doi:10.1090/S0002-9947-1944-0011087-2}.

\bibitem{Baker1975}
A.~Baker,
\newblock \emph{Transcendental Number Theory},
\newblock Cambridge University Press, 1975.

\bibitem{BakerWustholz1993}
A.~Baker and G.~W\"ustholz,
\newblock \emph{Logarithmic forms and group varieties},
\newblock Journal f\"ur die reine und angewandte Mathematik \textbf{442} (1993), 19--62.

\bibitem{Butler2017}
L.~A.~Butler,
\newblock \emph{A Diophantine approach to the three and four exponentials conjectures},
\newblock Ramanujan Journal \textbf{42} (2017), 199--221.

\bibitem{DasguptaKakde2026}
S.~Dasgupta and M.~Kakde,
\newblock \emph{Ranks of matrices of logarithms of algebraic numbers II: The matrix coefficient conjecture},
\newblock Advances in Mathematics \textbf{487} (2026), 110753.
\newblock \href{https://doi.org/10.1016/j.aim.2025.110753}{doi:10.1016/j.aim.2025.110753};
\href{https://arxiv.org/abs/2408.08178}{arXiv:2408.08178}.

\bibitem{CDTProducts}
F.~Calegari, V.~Dimitrov, and Y.~Tang,
\newblock \emph{The linear independence of $1$, $\zeta(2)$, and $L(2,\chi_{-3})$},
\newblock arXiv:2408.15403, 2024--2026; see Section~14 for products of two logarithms.

\bibitem{CDTHolonomy}
F.~Calegari, V.~Dimitrov, and Y.~Tang,
\newblock \emph{Arithmetic holonomy bounds and effective Diophantine approximation},
\newblock Proceedings of the International Congress of Mathematicians 2026, vol.~3,
356--376.
\newblock \href{https://doi.org/10.1137/25M1806776}{doi:10.1137/25M1806776};
\href{https://arxiv.org/abs/2510.04156}{arXiv:2510.04156}.

\bibitem{Kawashima2026}
M.~Kawashima,
\newblock \emph{Linear independence of periods related to polylogarithms},
\newblock arXiv:2605.17756v1, May 18, 2026.

\bibitem{Matveev2000}
E.~M.~Matveev,
\newblock \emph{An explicit lower bound for a homogeneous rational linear form in the logarithms of algebraic numbers. II},
\newblock Izvestiya: Mathematics \textbf{64} (2000), 1217--1269.

\bibitem{Waldschmidt2000}
M.~Waldschmidt,
\newblock \emph{Diophantine Approximation on Linear Algebraic Groups},
\newblock Springer, 2000.

\bibitem{Waldschmidt2023}
M.~Waldschmidt,
\newblock \emph{The Four Exponentials Problem and Schanuel's Conjecture},
\newblock in \emph{Mathematics Going Forward}, Lecture Notes in Mathematics, 2023, 579--592.

\bibitem{SuppliedReport}
The ProveIt project,
\newblock \emph{The Alaoglu--Erd\H{o}s Integer-Exponent Conjecture: Unified report and continuations},
\newblock supplied source file
\texttt{/Article/alaoglu-erdos-unified-report/alaoglu-erdos-unified-report.tex},
August 21, 2026.

\end{thebibliography}

\end{document}
